% GENERATED FILE. Edit canonical lessons, then run npm run generate:books.
% canonical-source-hash: fnv1a64:01b9b6cf8a92cfb8
% canonical-lessons: HI-C352-qaid, HI-C352-jaj, HI-C352-faisla, HI-C352-dosi, HI-C352-zamanat

\chapter{In Court}
\label{ch:hi-in-court}
\hlchaptermodality{352}

\begin{chapteropening}
\textbf{By the end of this chapter:} \emph{I can follow what happens in a court case.}

Complete the last lesson of chapter 352: I can follow what happens in a court case.
\end{chapteropening}

\section[\texorpdfstring{\hi{क़ैद} (qaid)}{क़ैद (qaid)}]{\hi{क़ैद} (qaid) — imprisonment, a prison sentence}

\label{lesson:HI-C352-qaid}



\begin{quote}
Before the new one: say the Hindi for proper, then the Hindi for justice.
\end{quote}

\subsection*{You'll want to know: \hi{क़ैद}}
\textbf{\hi{क़ैद}} — \emph{qaid} — \textquotedblleft{}imprisonment, a prison sentence\textquotedblright{}.

\begin{quote}
\hi{अदालत} \hi{ने} \hi{चोर} \hi{को} \hi{दो} \hi{साल} \hi{की} \hi{क़ैद} \hi{सुनाई।}
\end{quote}

\begin{quote}
\emph{adālat ne cor ko do sāl kī qaid sunāī.}
\end{quote}

\begin{quote}
The court sentenced the thief to two years in prison.
\end{quote}

Feminine.

\begin{cousinweb}
\textbf{\hi{क़ैद}} (\emph{qaid}) is usually traced to Arabic \emph{qayd}, \textquotedblleft{}a fetter, a bond\textquotedblright{}, which reached Hindi through Persian.
\end{cousinweb}

\begin{grammarlens}[title={What you've built}]
One more word for giving reasons and opinions.
\end{grammarlens}

\subsection*{Your turn}
\begin{itemize}
\raggedright
  \item \emph{Say it:} \emph{qaid}
  \item \emph{Say it:} \emph{qaid}, once more
  \item \emph{Say it:} \emph{adālat ne cor ko do sāl kī qaid sunāī.} — The court sentenced the thief to two years in prison.
  \item \emph{Recall:} say the Hindi for proper, then the Hindi for justice, then say \emph{qaid} again
\end{itemize}

\begin{tcolorbox}[breakable,colback=teal!4,colframe=teal!35!black,title={Before you move on}]
What does \hi{क़ैद} mean? (\textquotedblleft{}Imprisonment, a prison sentence\textquotedblright{}.) Say it once more.
\end{tcolorbox}
\section[\texorpdfstring{\hi{जज} (jaj)}{जज (jaj)}]{\hi{जज} (jaj) — a judge}

\label{lesson:HI-C352-jaj}



\begin{quote}
Before the new one: say the Hindi for justice, then the Hindi for imprisonment.
\end{quote}

\subsection*{You'll want to know: \hi{जज}}
\textbf{\hi{जज}} — \emph{jaj} — \textquotedblleft{}a judge\textquotedblright{}.

\begin{quote}
\hi{जज} \hi{ने} \hi{दोनों} \hi{पक्षों} \hi{की} \hi{बात} \hi{सुनी।}
\end{quote}

\begin{quote}
\emph{jaj ne donõ pakṣoṁ kī bāt sunī.}
\end{quote}

\begin{quote}
The judge heard both sides.
\end{quote}

Masculine or feminine, by the person meant.

\begin{cousinweb}
\textbf{\hi{जज}} (\emph{jaj}) is a loan from English \emph{judge}.
\end{cousinweb}

\begin{grammarlens}[title={What you've built}]
One more word for giving reasons and opinions.
\end{grammarlens}

\subsection*{Your turn}
\begin{itemize}
\raggedright
  \item \emph{Say it:} \emph{jaj}
  \item \emph{Say it:} \emph{jaj}, once more
  \item \emph{Say it:} \emph{jaj ne donõ pakṣoṁ kī bāt sunī.} — The judge heard both sides.
  \item \emph{Recall:} say the Hindi for justice, then the Hindi for imprisonment, then say \emph{jaj} again
\end{itemize}

\begin{tcolorbox}[breakable,colback=teal!4,colframe=teal!35!black,title={Before you move on}]
What does \hi{जज} mean? (\textquotedblleft{}A judge\textquotedblright{}.) Say it once more.
\end{tcolorbox}
\section[\texorpdfstring{\hi{फ़ैसला} (faislā)}{फ़ैसला (faislā)}]{\hi{फ़ैसला} (faislā) — a verdict, a decision}

\label{lesson:HI-C352-faisla}



\begin{quote}
Before the new one: say the Hindi for imprisonment, then the Hindi for a judge.
\end{quote}

\subsection*{You'll want to know: \hi{फ़ैसला}}
\textbf{\hi{फ़ैसला}} — \emph{faislā} — \textquotedblleft{}a verdict, a decision\textquotedblright{}.

\begin{quote}
\hi{अदालत} \hi{का} \hi{फ़ैसला} \hi{कल} \hi{आएगा।}
\end{quote}

\begin{quote}
\emph{adālat kā faislā kal āegā.}
\end{quote}

\begin{quote}
The court's verdict will come tomorrow.
\end{quote}

Masculine.

\begin{cousinweb}
\textbf{\hi{फ़ैसला}} (\emph{faislā}) is usually traced to Arabic \emph{faiṣala}, \textquotedblleft{}a decision\textquotedblright{}, which reached Hindi through Persian.
\end{cousinweb}

\begin{grammarlens}[title={What you've built}]
One more word for giving reasons and opinions.
\end{grammarlens}

\subsection*{Your turn}
\begin{itemize}
\raggedright
  \item \emph{Say it:} \emph{faislā}
  \item \emph{Say it:} \emph{faislā}, once more
  \item \emph{Say it:} \emph{adālat kā faislā kal āegā.} — The court's verdict will come tomorrow.
  \item \emph{Recall:} say the Hindi for imprisonment, then the Hindi for a judge, then say \emph{faislā} again
\end{itemize}

\begin{tcolorbox}[breakable,colback=teal!4,colframe=teal!35!black,title={Before you move on}]
What does \hi{फ़ैसला} mean? (\textquotedblleft{}A verdict, a decision\textquotedblright{}.) Say it once more.
\end{tcolorbox}
\section[\texorpdfstring{\hi{दोषी} (doṣī)}{दोषी (doṣī)}]{\hi{दोषी} (doṣī) — guilty}

\label{lesson:HI-C352-dosi}



\begin{quote}
Before the new one: say the Hindi for a judge, then the Hindi for a verdict.
\end{quote}

\subsection*{You'll want to know: \hi{दोषी}}
\textbf{\hi{दोषी}} — \emph{doṣī} — \textquotedblleft{}guilty\textquotedblright{}.

\begin{quote}
\hi{अदालत} \hi{ने} \hi{उसे} \hi{दोषी} \hi{माना।}
\end{quote}

\begin{quote}
\emph{adālat ne use doṣī mānā.}
\end{quote}

\begin{quote}
The court found him guilty.
\end{quote}

\begin{cousinweb}
\textbf{\hi{दोषी}} (\emph{doṣī}) is a tatsama, borrowed from Sanskrit \emph{doṣin}, \textquotedblleft{}guilty of an offence\textquotedblright{} (Monier-Williams).
\end{cousinweb}

\begin{grammarlens}[title={What you've built}]
One more word for giving reasons and opinions.
\end{grammarlens}

\subsection*{Your turn}
\begin{itemize}
\raggedright
  \item \emph{Say it:} \emph{doṣī}
  \item \emph{Say it:} \emph{doṣī}, once more
  \item \emph{Say it:} \emph{adālat ne use doṣī mānā.} — The court found him guilty.
  \item \emph{Recall:} say the Hindi for a judge, then the Hindi for a verdict, then say \emph{doṣī} again
\end{itemize}

\begin{tcolorbox}[breakable,colback=teal!4,colframe=teal!35!black,title={Before you move on}]
What does \hi{दोषी} mean? (\textquotedblleft{}Guilty\textquotedblright{}.) Say it once more.
\end{tcolorbox}
\section[\texorpdfstring{\hi{ज़मानत} (zamānat)}{ज़मानत (zamānat)}]{\hi{ज़मानत} (zamānat) — bail}

\label{lesson:HI-C352-zamanat}



\begin{quote}
Before the new one: say the Hindi for a verdict, then the Hindi for guilty.
\end{quote}

\subsection*{You'll want to know: \hi{ज़मानत}}
\textbf{\hi{ज़मानत}} — \emph{zamānat} — \textquotedblleft{}bail\textquotedblright{}.

\begin{quote}
\hi{जज} \hi{ने} \hi{उसे} \hi{ज़मानत} \hi{दी।}
\end{quote}

\begin{quote}
\emph{jaj ne use zamānat dī.}
\end{quote}

\begin{quote}
The judge granted him bail.
\end{quote}

Feminine.

\begin{cousinweb}
\textbf{\hi{ज़मानत}} (\emph{zamānat}) is usually traced to Arabic \emph{ḍamāna}, \textquotedblleft{}a guarantee\textquotedblright{}, which reached Hindi through Persian. Persian and Hindi say \emph{zamānat}.
\end{cousinweb}

\begin{grammarlens}[title={What you've built}]
One more word for giving reasons and opinions.
\end{grammarlens}

\subsection*{Your turn}
\begin{itemize}
\raggedright
  \item \emph{Say it:} \emph{zamānat}
  \item \emph{Say it:} \emph{zamānat}, once more
  \item \emph{Say it:} \emph{jaj ne use zamānat dī.} — The judge granted him bail.
  \item \emph{Recall:} say the Hindi for a verdict, then the Hindi for guilty, then say \emph{zamānat} again
\end{itemize}

\begin{tcolorbox}[breakable,colback=teal!4,colframe=teal!35!black,title={Before you move on}]
What does \hi{ज़मानत} mean? (\textquotedblleft{}Bail\textquotedblright{}.) Say it once more.
\end{tcolorbox}
